\begin{solution}{5.01}
According to Einstein's photoelectric equation, the maximum kinetic energy $E_{k\max}$ of the
photoelectrons is
\begin{equation*}
E_{k\max} = h\nu - \phi_0,
\end{equation*}
where $h\nu$ is the energy of the incident photon and $\phi_0$ is the work function of the
metal. Let the stopping voltage be $V_s$; then $eV_s = E_{k\max}$. Substituting this into the
above, the work function of the metal is found to be
\begin{equation*}
\phi_0 = h\nu - eV_s = (1240/200) - 2.6 = 3.6\unitv{eV}.
\end{equation*}

When ultraviolet light of wavelength $\lambda = 300\unitv{nm}$ is used instead, the stopping
voltage $V_s$ is
\begin{equation*}
V_s = (hc/\lambda - \phi_0)/e = 1240/300 - 3.6 = 0.53\unitv{V}.
\end{equation*}

If visible light is used to irradiate the surface of this metal, then, since the wavelength of
visible light is greater than $400\unitv{nm}$ and its energy is less than $3.1\unitv{eV}$,
while the minimum work necessary for an electron to escape from the metal is
$\phi_0 = 3.6\unitv{eV}$, no photoelectric effect can occur.
\end{solution}

\begin{solution}{5.02}
\textbf{Method 1.}

For part (h), if the production of a large number of protons is regarded as the result of
$\gamma$ photons striking the hydrogen nuclei in paraffin (this phenomenon is usually called
the Compton effect of the proton), then, analogously to the Compton effect, we may regard the
process as an elastic collision between a $\gamma$ photon of energy $h\nu$ and a free proton
initially at rest. Since a photon is involved and the recoil velocity of the proton may
approach the speed of light, relativistic expressions must be used.

From the conservation of energy and mass,
\begin{equation*}
\frac{hc}{\lambda} + M_pc^2 = \frac{hc}{\lambda'} + \sqrt{(pc)^2 + (M_pc^2)^2}.
\end{equation*}

From the conservation of momentum,
\begin{equation*}
p^2 = \left(\frac{h}{\lambda}\right)^2 + \left(\frac{h}{\lambda'}\right)^2
      - 2\,\frac{h^2}{\lambda\lambda'}\cos\theta,
\end{equation*}
where $\lambda$ and $\lambda'$ are respectively the wavelengths of the $\gamma$ photon before
and after the collision; $M_pc^2$ is the rest energy of the proton,
$M_pc^2 = 0.511\times 1836 = 938\unitv{MeV}$; $p$ is the recoil momentum of the proton; and
$\theta$ is the scattering angle of the photon. Combining the above two equations and
eliminating $p$, we obtain
\begin{equation*}
\Delta\lambda = \lambda' - \lambda = \frac{hc}{M_pc^2}(1 - \cos\theta).
\end{equation*}

When $\theta = \pi$, the recoil energy of the proton is maximum, namely
\begin{equation*}
E_k = \left(\frac{hc}{\lambda}\right) - \left(\frac{hc}{\lambda'}\right)
    = \left(\frac{hc}{\lambda}\right)
      - \frac{hc}{\lambda + \dfrac{2hc}{M_pc^2}}
    = 5.7\unitv{MeV}.
\end{equation*}

From the above equation the energy $h\nu$ of the incident photon can be solved:
\begin{align*}
h\nu = \frac{hc}{\lambda}
     &= \frac{1}{2}\left(E_k + \sqrt{E_k^2 + 2E_kM_pc^2}\right)\\
     &= \frac{1}{2}\times\left(5.7 + \sqrt{5.7^2 + 2\times 5.7\times 938}\right)
      = 55\unitv{MeV}.
\end{align*}

That is, to produce, after scattering, a proton with the maximum recoil energy of
$5.7\unitv{MeV}$, the incident $\gamma$ photon must have an energy of $55\unitv{MeV}$. This
result is inconsistent with the result of (a); therefore this particle is not a $\gamma$
photon, unless the conservation of energy and the conservation of momentum are abandoned.

\textbf{Method 2.}

Analogously to Method 1, we obtain
\begin{equation*}
\Delta\lambda = \lambda' - \lambda = \frac{hc}{M_pc^2}(1 - \cos\theta).
\end{equation*}
When $\theta = \pi$, the recoil energy of the proton is maximum, namely
\begin{equation*}
E_k = \left(\frac{hc}{\lambda}\right) - \left(\frac{hc}{\lambda'}\right)
    = \left(\frac{hc}{\lambda}\right)
      - \frac{hc}{\lambda + \dfrac{2hc}{M_pc^2}}
    = h\nu - \frac{h\nu}{1 + \dfrac{2h\nu}{M_pc^2}}.
\end{equation*}
From part (a) we know that $h\nu = 15\unitv{MeV}$. Substituting this into the above equation
gives
\begin{equation*}
E_k = 15 - 15/(1 + 2\times 15/938) = 0.46\unitv{MeV}.
\end{equation*}
This result is inconsistent with the result of (b); therefore this particle is not a $\gamma$
photon, unless the conservation of energy and the conservation of momentum are abandoned.

\begin{commentary}
The problem is based on some experiments connected with the discovery of the neutron. In 1920,
after much reflection, Rutherford put forward his famous neutron hypothesis:\footnote{E.~Rutherford,
\emph{Proc. Roy. Soc.} A97 (1920) p.~374.} he considered it possible that there exists a new
kind of neutral atom, perhaps a double particle formed by an electron and a hydrogen nucleus
bound together more tightly. He imagined that this atom would have many novel properties,
properties that are necessary for explaining the composition of the nuclei of heavy elements.

In order to test Rutherford's hypothesis, scientists carried out a large amount of
experimental work. In 1930 the German W. Bothe and his student H. Becker were the first to
publish the results of experiments in which particles produced by radioactive polonium
bombarded beryllium: when $\alpha$ particles strike beryllium, the beryllium emits what was
then regarded as a more penetrating kind of $\gamma$ radiation. In 1932 the Joliot-Curies
repeated these experiments with a very strong polonium $\alpha$ source and measured the
penetrating power of this unknown radiation through lead. They found that the intensity of
the radiation emitted by beryllium decreased by only one half in passing through $4.7\unitv{cm}$
of lead. If this radiation were $\gamma$ rays, its energy would be about $15\unitv{MeV}$. This
value is much higher than the energy of the $\gamma$ rays emitted by natural radioactive
substances. They also found that this unknown radiation, on falling on hydrogen-containing
paraffin, produces many protons of very high speed. Unfortunately, they did not carefully
check whether, according to the conservation of momentum and the conservation of energy,
$\gamma$ particles of energy $15\unitv{MeV}$ striking hydrogen nuclei could produce such a
high-speed proton stream. Like Bothe, they regarded this radiation emitted by beryllium as
$\gamma$ rays, and considered the production of the proton stream to be the result of
$\gamma$ particles striking hydrogen nuclei, a kind of special phenomenon similar to the
Compton effect. In this way both Bothe and the Joliot-Curies missed the opportunity of
discovering the neutron.

Rutherford's early student and capable assistant J. Chadwick had been engaged in experiments
searching for the neutron ever since 1921. In 1932, when he read the report of the
Joliot-Curies on the astonishing properties of the beryllium radiation, he told Rutherford of
the views of the Joliot-Curies. Rutherford did not believe that this was the Compton effect,
and thought it very likely that the neutron, sought for many years, had appeared here.

Chadwick quickly rechecked the experimental results of Bothe and of the Joliot-Curies. By
measuring the range of these protons in air, he found the energy of the protons to be
$5.7\unitv{MeV}$. He also extended the experiment, using the beryllium radiation to irradiate
many other substances, including the gases hydrogen, helium, nitrogen and argon, and
determined the energy of the recoil atoms produced thereby, from which he calculated the
energy of the incident $\gamma$ rays. Chadwick found that this result depended on the gas
used. For example, for recoil nitrogen atoms of energy $1.4\unitv{MeV}$, the energy of the
incident $\gamma$ rays should be $90\unitv{MeV}$, while to produce recoil protons of energy
$5.7\unitv{MeV}$ the energy of the $\gamma$ rays should be $55\unitv{MeV}$. If the beryllium
radiation consisted of neutral particles of mass approximately equal to the proton mass, then
when these particles (called neutrons by Chadwick, following Rutherford) make a head-on
collision with the atoms of a gas, the maximum recoil energy will appear. As for the recoil
protons, since the proton and the neutron have equal masses, in a head-on collision the
proton will acquire the whole kinetic energy of the neutron; hence the energy of the neutrons
emitted from beryllium equals the maximum recoil energy of the protons, $5.7\unitv{MeV}$.
Neutrons of this energy, incident on nitrogen gas, can produce recoil nitrogen atoms of
energy $1.4\unitv{MeV}$.

In the same year Chadwick published a paper entitled ``The Existence of the Neutron'',\footnote{J.~Chadwick,
\emph{Proc. Roy. Soc.} (London) A136 (1932) p.~692.} in which he described his experimental
results and theoretical analysis in detail. He pointed out in the paper that only by
abandoning the conservation of energy and momentum could the hypothesis that the beryllium
radiation consists of $\gamma$ rays be supported. The neutron hypothesis, on the other hand,
gives a direct and simple explanation of the experimental facts. It is self-consistent and
offers new insight into the structure of atomic nuclei. Using the nuclear reaction
${}^{4}\mathrm{He} + {}^{11}\mathrm{B} \rightarrow {}^{14}\mathrm{N} + \mathrm{n}$, Chadwick
concluded that the mass of the neutron lies roughly between 1.005 and 1.008.
\end{commentary}
\end{solution}

\begin{solution}{5.03}
(1) From the Compton shift formula
\begin{equation*}
\Delta\lambda = \lambda_c(1 - \cos\theta),
\end{equation*}
where $\lambda_c = h/(m_ec) = 0.0243\unitv{\mathring{A}}$ is the Compton wavelength of the
electron (for the derivation of the formula see Solution 5.02), we obtain:

(a) If after scattering the electron moves along the direction of incidence of the photon,
then $\theta = \pi$, and
\begin{equation*}
\Delta\lambda = 2\lambda_c = 0.00486\unitv{nm}.
\end{equation*}

(b) If the photon still moves along the original direction of incidence, then $\theta = 0$ and
$\Delta\lambda = 0$.

\begin{figbox}{5.03}
\begin{tikzpicture}[scale=1.0]
  \coordinate (O) at (0,0);
  \coordinate (A) at (5,0);
  \coordinate (B) at (5,2.1);
  \draw[force] (O) -- (A) node[midway,below] {$h/\lambda$};
  \draw[force] (O) -- (B) node[midway,above left] {$h/\lambda'$};
  \draw[force] (B) -- (A) node[midway,right] {$p$};
  \draw (1.0,0) arc (0:22.8:1.0);
  \node at (1.28,0.28) {$\theta$};
\end{tikzpicture}
\end{figbox}

(2)(a) From the conservation of energy and mass,
\begin{equation*}
\frac{hc}{\lambda} + \sqrt{p_e^2c^2 + E_0^2} = \frac{hc}{\lambda'} + E_0.
\end{equation*}

From the conservation of momentum,
\begin{equation*}
p_e^2 = \left(\frac{h}{\lambda}\right)^2 + \left(\frac{h}{\lambda'}\right)^2
      - 2\,\frac{h^2}{\lambda\lambda'}\cos\theta,
\end{equation*}
where $\lambda$ and $\lambda'$ are respectively the wavelengths of the photon before and after
the collision; $E_0$ is the rest energy of the electron; $p_e$ is the initial momentum of the
electron; and $\theta$ is the scattering angle of the photon. Combining the above equations
and eliminating $p_e$, we obtain
\begin{align*}
\Delta\lambda = \lambda - \lambda' &= \lambda_c(1 - \cos\theta)\\
&= \lambda_c(1 - \cos 60\degree) = 0.0012\unitv{nm}.
\end{align*}

(b) It is known that the energy of the incident photon equals the rest energy of the electron;
hence the wavelength of the incident photon is $\lambda = \lambda_c$. From (a) we have
$\Delta\lambda = \lambda_c/2$. Substituting $\lambda$ and $\Delta\lambda$ into the following
expression, we obtain the energy $h\nu'$ of the scattered photon:
\begin{equation*}
h\nu' = \frac{hc}{\lambda'} = \frac{hc}{\lambda - \Delta\lambda} = 2E_0 = 1.02\unitv{MeV}.
\end{equation*}
From the conservation of energy, the kinetic energy $E_k$ of the electron before the collision
is found to be
\begin{equation*}
E_k = \left(\frac{hc}{\lambda'}\right) - \left(\frac{hc}{\lambda}\right)
    = E_0 = 0.511\unitv{MeV}.
\end{equation*}

\begin{commentary}
Part (1) of this problem discusses the Compton effect: a high-frequency $\gamma$ photon makes
an elastic collision with a free electron that is approximately at rest; the photon transfers
part of its energy to the electron, this becoming the recoil kinetic energy of the electron,
while the wavelength of the scattered photon becomes longer. Part (2) discusses the reverse
process of the Compton effect: the collision of a high-energy relativistic electron with a
low-frequency photon, after which the energy (frequency) of the scattered photon becomes
larger and its wavelength shorter; it is therefore called the inverse Compton effect. The
inverse Compton effect is an important non-thermal radiation mechanism for the production of
$\gamma$ rays.

The Compton effect was discovered by A.~H. Compton in 1923, who also gave its theoretical
explanation. Following the photoelectric effect, it provided further confirmation of the
particle nature of radiation; hence it is of extremely profound significance both
theoretically and experimentally. For his series of experiments on the Compton effect and his
theoretical explanation of it, Compton shared the 1927 Nobel Prize in Physics with Wilson of
Britain.\footnote{The Chinese physicist Wu Youxun made outstanding contributions to Compton
scattering experiments; see Guan Weiyan, ``A brief account of Professor Wu Youxun'',
\emph{Physics}, 11 (1982) 457.} During the ten years from 1930 to 1940 Compton devoted himself
to the study of cosmic rays and discovered the inverse Compton effect. The Compton effect and
the inverse Compton effect are of great importance in astrophysics, because the interaction
of $\gamma$ photons and electrons in cosmic rays is an important aspect of the Compton effect,
while the inverse Compton effect of relativistic electrons is an important topic of
contemporary astrophysics. For example, in April 1991, following the Hubble telescope, NASA
launched the second giant space observatory, the Compton Gamma Ray Observatory, which can
image cosmic $\gamma$-ray sources and their background and measure their spectra and their
time variation. The Compton imaging telescope with which it is equipped uses the Compton
effect to detect $\gamma$ rays of energy $1$--$30\unitv{MeV}$ in the universe. As another
example, the large number of high-temperature electrons present in the ionised gas of clusters
of galaxies produce inverse Compton scattering of the photons of the cosmic microwave
background, giving rise to the so-called Sunyaev--Zel'dovich (SZ) effect. The SZ effect is a
very effective means of studying clusters of galaxies; in particular, the SZ effect depends
only on the intrinsic properties of the cluster and not on its position, which will be
favourable for discovering clusters of high redshift. It is therefore of extremely great
significance for studying the evolution of clusters of galaxies. Astronomical telescopes
built on the basis of the SZ effect will help cosmologists to find clusters of galaxies that
are too distant or too faint to be discovered by other methods. These observations in turn
help cosmologists to discover how clusters evolve, and may make it possible to locate the
dark matter in the universe. The SZ cluster survey telescope now under construction will use
the SZ effect and provide an entirely new and effective way of determining cosmological
parameters.
\end{commentary}
\end{solution}

\begin{solution}{5.04}
(1) As shown in Fig.~5.04(a), before the collision the energy of the electron is $E$, its
momentum is $p_x = p$, $p_y = 0$; the energy of the photon is $h\nu$, its momentum is
$p_x = -h\nu/c$, $p_y = 0$. As shown in Fig.~5.04(b), after the collision the momentum of the
electron is $p'$, directed at an angle to the $x$-direction; the energy of the photon is
$h\nu'$, its momentum is $h\nu'/c$, making an angle $\theta$ with the $x$-direction.

\begin{figbox}{5.04(a)}
\begin{tikzpicture}[scale=0.95]
  \node at (-2.9,0.55) {electron};
  \node at (-2.9,-0.55) {$(E,p,0)$};
  \node at (2.9,0.55) {photon};
  \node at (2.9,-0.55) {$(h\nu,-h\nu/c,0)$};
  \draw[force] (-1.75,0) -- (-0.2,0);
  \draw[force] (1.75,0) -- (0.2,0);
\end{tikzpicture}
\end{figbox}

\begin{figbox}{5.04(b)}
\begin{tikzpicture}[scale=1.0]
  \coordinate (O) at (0,0);
  \coordinate (A) at (5.6,0);
  \coordinate (B) at (4.3,2.1);
  \draw[force] (O) -- (A) node[midway,below] {$p-h\nu/c$};
  \draw[force] (O) -- (B) node[midway,above left] {$h\nu'/c$};
  \draw[force] (B) -- (A) node[midway,above right] {$p'$};
  \draw (1.1,0) arc (0:26:1.1);
  \node at (1.38,0.30) {$\theta$};
  \node[above left] at (B) {scattered photon};
  \node[below] at (4.55,0.45) {scattered electron};
\end{tikzpicture}
\end{figbox}

From the conservation of energy,
\begin{equation}
E + h\nu = \sqrt{p'^2c^2 + E_0^2} + h\nu', \tag{1}
\end{equation}
where $E_0$ is the rest energy of the electron. From the conservation of momentum,
\begin{equation}
p'^2c^2 = (h\nu')^2 + (pc - h\nu)^2 - 2(h\nu')(pc - h\nu)\cos\theta. \tag{2}
\end{equation}

Substituting (2) into (1), we obtain
\begin{equation*}
h\nu' = \frac{(E + pc)h\nu}{E + h\nu + h\nu\cos\theta - pc\cos\theta}.
\end{equation*}

Since
\begin{equation*}
E = \sqrt{p^2c^2 + E_0^2} = 20\unitv{GeV} \gg E_0,
\end{equation*}
we have $E \simeq pc$, and hence
\begin{equation*}
h\nu' \simeq \frac{2h\nu E}{h\nu(1 + \cos\theta) + E - pc\cos\theta}.
\end{equation*}

(2) When $\theta = 0$, the energy of the scattered photon is maximum, its value being
\begin{equation*}
(h\nu')_{\max} = \frac{2h\nu E}{2h\nu + E - pc}.
\end{equation*}

In the extreme relativistic case,
\begin{equation*}
E = \sqrt{p^2c^2 + E_0^2} \simeq pc + \frac{E_0^2}{2pc},
\end{equation*}
\begin{equation*}
E - pc \simeq \frac{E_0^2}{2E} \simeq \frac{(0.5\times 10^6)^2}{2\times 2\times 10^{10}}
\approx 6\unitv{eV}.
\end{equation*}

For a typical visible photon,
\begin{equation*}
h\nu = \frac{hc}{\lambda} \approx \frac{12\,400}{5\,000} \approx 2.5\unitv{eV}.
\end{equation*}
Thus the maximum photon energy obtainable by this technique is
\begin{equation*}
(h\nu')_{\max} = \frac{2\times 2.5\times 20}{2\times 2.5 + 6} = 9\unitv{GeV}.
\end{equation*}
\end{solution}

\begin{solution}{5.05}
Suppose that the electron moves along the $y$-direction with momentum $p_y$. From the
relativistic momentum--energy relation,
\begin{equation*}
cp_y = \sqrt{E^2 - E_0^2},
\end{equation*}
where $c$ is the speed of light, $E$ is the energy of the electron, $E = 1\unitv{MeV}$, and
$E_0$ is the rest energy of the electron, $E_0 = 0.511\unitv{MeV}$.

When a fast-moving electron travels through a cloud chamber, the ions produced along its path
form condensation centres, causing the water vapour to condense into small droplets. The
linear size of a droplet is about $10^{-5}\unitv{m}$; therefore in this case
$\Delta x = 10^{-5}\unitv{m}$. From the uncertainty relation we may assume that the
uncertainty in the momentum of the electron in the $x$-direction is about
\begin{equation*}
\Delta p_x = \hbar/\Delta x,
\end{equation*}
where $\hbar = h/2\pi = 1973\unitv{eV.\mathring{A}}$.

Consequently the motion of the electron will deviate from the $y$-direction by an angle
$\theta$, and the uncertainty $\Delta\theta$ of this angle is about
\begin{align*}
\Delta\theta = \frac{\Delta p_x}{p_y}
&= \frac{\hbar c}{\Delta x\sqrt{E^2 - E_0^2}}\\
&= \frac{1973}{10^{-5}\times 10^{10}\times\sqrt{10^{12} - (0.511\times 10^6)^2}}\\
&= 2.3\times 10^{-8},
\end{align*}
\begin{equation*}
\theta = 2.3\times 10^{-8}\unitv{rad}.
\end{equation*}
It is thus seen that the track of the electron will not differ noticeably from a straight
line.

\begin{figbox}{5.05}
\begin{tikzpicture}[scale=1.0]
  \draw[axis] (0,0) -- (0,4.0) node[right] {$p_y$};
  \draw[axis] (0,0) -- (2.6,0) node[below] {$\Delta p_x$};
  \draw[force] (0,0) -- (2.3,3.5) node[right] {$p$};
  \draw[dashedline] (2.3,0) -- (2.3,3.5);
  \draw[dashedline] (0,3.5) -- (2.3,3.5);
  \draw (0,1.2) arc (90:56.7:1.2);
  \node at (0.38,1.48) {$\theta$};
\end{tikzpicture}
\end{figbox}

\begin{commentary}
(1) The estimate of $\Delta p_x$ in the solution can evidently be applied to the momentum
components in every direction; here, in order to estimate the maximum deviation of the
electron track from a straight line, only the transverse momentum uncertainty is considered.

(2) The particles existing in nature have wave properties. When a particle moves with a
definite momentum $p$, its wavelength is $\lambda = h/p$. This relation between wavelength and
momentum holds for all real particles. Under certain conditions, however, we may still regard
a particle as a classical particle; the uncertainty relation helps us to estimate the range of
validity of the classical description of particles. For example, under the conditions of this
problem the energy of the electron is very high and the region to which it is confined is
relatively large, so that $\Delta p/p \ll 1$; hence the electron may be regarded as a classical
particle moving along a straight track. At this time the de Broglie wavelength of the electron
is very short, $\lambda = h/p = 1.4\times 10^{-12}\unitv{m} \ll 10^{-5}\unitv{m}$, which is
similar to the condition that makes geometrical optics valid in optics. If an electron with a
kinetic energy of a few electron-volts were confined in a space of atomic size, then in this
case $\Delta p \sim p$ for the electron, and the electron would have no definite track. The
classical description of the particle would no longer be applicable.

(3) The $p$ and $E$ in the above formulae are both relativistic, which also shows that there
is a close relation between quantum theory and relativity. Only when the kinetic energy $E_k$
of the particle is much smaller than its rest energy $E_0$ may the non-relativistic
momentum--energy relation $pc = \sqrt{2E_0E_k}$ be used.
\end{commentary}
\end{solution}

\begin{solution}{5.06}
Let the distance of the electron from the nucleus be $r$, and take the order of magnitude of
the position uncertainty to be $\Delta x = r$. From the uncertainty relation
$\Delta x\,\Delta p \ge \hbar$ we have $(\Delta p)^2 = p^2 \ge \hbar^2/r^2$; hence

(1)
\begin{equation*}
E = \frac{\hbar^2}{2mr^2} - \frac{Ze^2}{4\pi\varepsilon_0r}.
\end{equation*}
When $E$ takes its minimum value the hydrogen-like ion is in its ground state. Let its radius
be $r_0$ and its energy be $E_1$; from $\left(\dfrac{\dd E}{\dd r}\right)_{r=r_0} = 0$ we obtain
$r_0$ and $E_1$ as
\begin{equation*}
r_0 = \frac{4\pi\varepsilon_0\hbar^2}{Ze^2m},
\end{equation*}
\begin{equation*}
E_1 = -\frac{mZ^2e^4}{2(4\pi\varepsilon_0\hbar)^2}.
\end{equation*}

(2)
\begin{equation*}
E = \sqrt{p^2c^2 + m^2c^4} - \frac{Ze^2p}{4\pi\varepsilon_0\hbar}.
\end{equation*}
Let the momentum of the electron when the hydrogen-like ion is in its ground state be $p_0$,
and its energy be $E_1'$. From $(\dd E/\dd p)_{p=p_0} = 0$ we obtain $p_0$ and $E_1'$ as
\begin{equation*}
p_0 = \frac{mcZe^2}{\sqrt{(4\pi\varepsilon_0\hbar c)^2 - (Ze^2)^2}},
\end{equation*}
\begin{equation*}
E_1' = mc^2\sqrt{1 - \left(\frac{Ze^2}{4\pi\varepsilon_0\hbar c}\right)^2}.
\end{equation*}

(3) When $\dfrac{Ze^2}{4\pi\varepsilon_0\hbar c} \ll 1$, that is, $Z \ll \alpha^{-1} = 137$,
expanding the right-hand side of the above expression in series and neglecting the
higher-order infinitesimals, the above expression can be written as
\begin{equation*}
E_1' = mc^2\left[1 - \frac{1}{2}\left(\frac{Ze^2}{4\pi\varepsilon_0\hbar c}\right)^2\right]
     = E_1 + mc^2.
\end{equation*}

$E_1$ and $E_1'$ differ by a constant $mc^2$; this is caused by the different choices of the
zero of energy in the two cases. Therefore $E_1$ is the classical limit of $E_1'$. When $Z$ is
small the non-relativistic approximation may be used.

\begin{commentary}
(1) The $r_0$ and $E_1$ obtained in part (1) agree exactly with the results of Bohr's theory.
This is because $\Delta x = r$ was wisely chosen rather than $2r$ or $r/2$; such a choice is
precisely reasonable. Of course, any reasonable choice of $\Delta x$ will give the correct
order of magnitude of the atomic size and of the atomic energy.

(2) In part (3), if one more term is retained in the series expansion of $E_1'$, one obtains
\begin{equation*}
E_1' = mc^2\left[1 - \frac{1}{2}\left(\frac{Ze^2}{4\pi\varepsilon_0\hbar c}\right)^2
      - \frac{1}{8}\left(\frac{Ze^2}{4\pi\varepsilon_0\hbar c}\right)^4\right].
\end{equation*}
To make the zero of energy agree with the non-relativistic case, we subtract $mc^2$ from the
above expression and thus obtain
\begin{equation*}
E_1' = -mc^2\left[\frac{1}{2}\left(\frac{Ze^2}{4\pi\varepsilon_0\hbar c}\right)^2
      + \frac{1}{8}\left(\frac{Ze^2}{4\pi\varepsilon_0\hbar c}\right)^4\right]
     = E_1 - \frac{mc^2}{8}\left(\frac{Ze^2}{4\pi\varepsilon_0\hbar c}\right)^4.
\end{equation*}
The term $-\dfrac{mc^2}{8}\left(\dfrac{Ze^2}{4\pi\varepsilon_0\hbar c}\right)^4
= -|E_1|\dfrac{Z^2\alpha^2}{4} = \Delta E_r$ is the relativistic correction to the kinetic
energy of the ground state of the hydrogen-like ion. The result for this term derived
rigorously in quantum mechanics is
\begin{equation*}
\Delta E_r = -|E_1|\,\frac{5Z^2\alpha^2}{4}.
\end{equation*}
For the hydrogen atom, $Z = 1$ and $\Delta E_r = -6.66\times 10^{-5}|E_1| \ll |E_1|$; it is
therefore called the fine-structure correction of the energy level. But for hydrogen-like
ions with $Z = 50$, this correction already exceeds 16\%.
\end{commentary}
\end{solution}

\begin{solution}{5.07}
(1) The circumference $S_n$ of the $n$-th orbit of the hydrogen atom and the de Broglie
wavelength $\lambda_n$ of the electron moving in that orbit are respectively
\begin{equation*}
S_n = 2\pi r_n, \qquad \lambda_n = \frac{h}{mv_n},
\end{equation*}
where $r_n$ and $v_n$ are respectively the radius of the $n$-th orbit and the orbital speed of
the electron, and $m$ is the mass of the electron. From the quantisation condition for the
angular momentum of the electron, $mv_nr_n = n\hbar$, we obtain
\begin{equation*}
\frac{h}{\lambda_n}\cdot 2\pi r_n = n\hbar.
\end{equation*}
Finally,
\begin{equation*}
\frac{S_n}{\lambda_n} = n.
\end{equation*}

(2) Substituting $n = 1$ into the above expression gives
\begin{equation*}
\lambda_1 = S_1/1 = 2\pi a_0,
\end{equation*}
where $a_0 = 0.0529\unitv{nm}$ is the Bohr radius. Substituting this value into the above
expression gives
\begin{equation*}
\lambda_1 = 0.332\unitv{nm}.
\end{equation*}

\begin{commentary}
The solution of part (1) shows that in the hydrogen atom only those orbits whose
circumference is an integral multiple of the de Broglie wavelength of the electron are
allowed. One may imagine that, just as a standing wave is formed on a string fixed at both
ends, the de Broglie wave of the electron is here a standing wave on a circle. Bohr's
quantisation condition for angular momentum is precisely the necessary condition for the
formation of such a standing wave. In 1924 de Broglie gave a physical explanation of Bohr's
quantisation condition for angular momentum; he pointed out that Bohr's quantisation
condition for angular momentum is equivalent to the standing-wave condition.

In fact, according to the uncertainty principle the orbit of an electron cannot be determined
precisely; moreover, the orbital angular momentum of the lowest energy state of the hydrogen
atom is zero, and not the value $\hbar$ predicted by Bohr. Later, Bohr's model of the hydrogen
atom was replaced by the quantum-mechanical theory of the hydrogen atom. Nevertheless, in
quantum mechanics the orbital angular momentum is still quantised.
\end{commentary}
\end{solution}

\begin{solution}{5.08}
The wave numbers of the spectral lines of the hydrogen atom and of the helium ion are
respectively given by the following formulae:
\begin{equation}
\tilde{\nu}_{\mathrm{H}} = R_{\mathrm{H}}\left(\frac{1}{m^2} - \frac{1}{n^2}\right), \tag{1}
\end{equation}
\begin{equation}
\tilde{\nu}_{\mathrm{He}} = R_{\mathrm{He}}Z^2\left(\frac{1}{m^2} - \frac{1}{n^2}\right)
= R_{\mathrm{He}}\left[\frac{1}{(m/2)^2} - \frac{1}{(n/2)^2}\right], \tag{2}
\end{equation}
where $m = 1, 2, 3, \ldots$; $n = m+1, m+2, m+3, \ldots$; and for the helium ion $Z = 2$.
Since the mass of the helium nucleus is greater than that of the hydrogen nucleus, the
Rydberg constant $R_{\mathrm{H}} \simeq R_{\mathrm{He}}$, but $R_{\mathrm{He}}$ is slightly
greater than $R_{\mathrm{H}}$.

Comparing (1) and (2), it is seen that when the quantum numbers $m$ and $n$ of the helium ion
and the quantum numbers of the hydrogen atom satisfy the conditions $(m/2)_{\mathrm{He}} =
(m)_{\mathrm{H}}$ and $(n/2)_{\mathrm{He}} = (n)_{\mathrm{H}}$, we have
$\tilde{\nu}_{\mathrm{H}} \simeq \tilde{\nu}_{\mathrm{He}}$, but
$\tilde{\nu}_{\mathrm{H}}$ is slightly smaller than $\tilde{\nu}_{\mathrm{He}}$. Therefore
some almost coincident doublets will appear in the spectrogram: the line of slightly smaller
wave number in each doublet is the spectral line emitted by the hydrogen atom; apart from
these, all the remaining lines belong to the spectrum of the helium ion. Since the intensity
and the spacing of the spectral lines both decrease in the direction of increasing wave
number, the magnitude of the wave number can be judged from their trend of variation.

\begin{commentary}
In 1897 the astronomer W.~H. Pickering discovered in the spectrum of the star $\xi$ in the
constellation Puppis a series of spectral lines very similar to the Balmer series of the
hydrogen atom; it is called the Pickering series. The figure below shows the first five lines
$H_\alpha$, $H_\beta$, $H_\gamma$, $H_\delta$ and $H_\varepsilon$ of the Pickering series and
of the Balmer series of the hydrogen atom. From the figure it can be seen that every other
line of the Pickering series almost coincides with a line of the Balmer series, but the
wavelengths of the two still differ slightly. There are other lines of the Pickering series
lying between adjacent lines of the Balmer series. Rydberg pointed out that the Pickering
series can be represented by
\begin{equation*}
\tilde{\nu}_P = R\left(\frac{1}{2^2} - \frac{1}{k^2}\right),
\qquad k = 2.5,\ 3,\ 3.5,\ 4,\ 4.5,\ \ldots
\end{equation*}
where $\tilde{\nu}_P$ is the wave number of the spectral line. The Pickering series cannot be
observed in hydrogen on the earth. At first people thought that it was emitted by a special
kind of hydrogen in the celestial bodies; later it was found that the Pickering series
appears when some helium is mixed with hydrogen, and only then was it made clear that the
Pickering series is emitted by the helium ion $\mathrm{He}^+$. In the solution of this
problem it is only necessary to set $m = 4$, $n = 5, 6, 7, \ldots$ in equation (2) to obtain
the expression for the Pickering series.

\begin{figbox}{5.08}
\begin{tikzpicture}[scale=0.92]
  \draw[thickline] (0,0) -- (10.5,0);
  \foreach \x/\lab in {1.0/25\,000, 5.2/20\,000, 9.0/15\,000}{
    \draw[thickline] (\x,0) -- (\x,-0.10);
    \node[below,font=\small] at (\x,-0.12) {$\lab$};
  }
  \node[below,font=\small] at (10.3,-0.12) {$\mathrm{cm}^{-1}$};
  % Balmer lines
  \foreach \x/\lab/\h in {2.05/{$H_\varepsilon$}/1.55, 2.95/{$H_\delta$}/1.35,
                          4.05/{$H_\gamma$}/1.45, 5.55/{$H_\beta$}/1.55,
                          8.15/{$H_\alpha$}/1.75}{
    \draw[line width=1.0pt] (\x,0) -- (\x,\h);
    \node[above,font=\small] at (\x,\h) {\lab};
  }
  % Pickering lines (weaker, interleaved)
  \foreach \x/\h in {1.55/0.60, 2.50/0.55, 3.50/0.65, 4.80/0.55,
                    6.60/0.50, 7.35/0.55}{
    \draw[line width=0.5pt] (\x,0) -- (\x,\h);
  }
\end{tikzpicture}
\end{figbox}
\end{commentary}
\end{solution}

\begin{solution}{5.09}
(1) Since the masses of the $\mu^-$ meson and the $\pi^+$ meson are of the same order of
magnitude, both mesons move about their common centre of mass under the Coulomb force exerted
by the other. Introducing the reduced mass
\begin{equation*}
\mu = \frac{m_\mu m_\pi}{m_\mu + m_\pi},
\end{equation*}
the lowest energy $E_1$ of the atom is
\begin{equation*}
E_1 = -\frac{1}{2}\mu(\alpha c)^2,
\end{equation*}
where $\alpha$ is the fine-structure constant and $c$ is the speed of light. Substituting
$m_\mu$ and $m_\pi$ into the above expression gives
\begin{equation*}
E_1 = -\frac{1}{2}(\alpha c)^2m_e\times\frac{106\times 140}{0.511\times(106 + 140)}
    = -13.6\times 118\unitv{eV} = -1.60\unitv{keV}.
\end{equation*}

(2) In the ground state the distance $r_1$ between the two mesons is
\begin{equation*}
r_1 = \frac{a_0m_e}{\mu} = \frac{0.529\unitv{\mathring{A}}}{118}
    = 4.48\times 10^{-4}\unitv{nm}.
\end{equation*}

(3) In the ground state the relative speed of the two mesons is $v_1 = \alpha c$. Neglecting
the overall motion of the atom, we have
\begin{equation*}
m_\pi v_\pi = m_\mu v_\mu,
\end{equation*}
\begin{equation*}
v_1 = v_\pi + v_\mu.
\end{equation*}
From these the speed of the $\mu^-$ meson is found to be
\begin{equation*}
v_\mu = \frac{\alpha cm_\pi}{m_\mu + m_\pi}
     = \frac{3\times 10^8\times 140}{137\times(106 + 140)}
     = 1.25\times 10^6\unitv{m/s}.
\end{equation*}

\begin{commentary}
An ordinary atom is composed of electrons, protons and neutrons. A so-called exotic atom is
an atom-like system formed, through electromagnetic interaction, when particles such as
$\mu^\pm$ mesons, $\tau^\pm$ mesons, $\pi^\pm$ mesons, $K^\pm$ mesons, positrons,
antiprotons, $\Sigma^\pm$ hyperons and $\Omega^-$ hyperons replace respectively the electrons
or the nucleus of an ordinary atom, or replace both. For example, a $\mu^-$ meson replacing an
electron of an atom forms a muonic atom, and a positron replacing the nucleus forms
positronium. The average lifetime of an exotic atom is not long, but there is enough time for
the exotic atom to be formed and detected before the replacing particle itself decays. The
best known of these is the antihydrogen atom, composed of an antiproton and a positron.
\end{commentary}
\end{solution}

\begin{solution}{5.10}
According to Gauss's law in electrostatics it can be shown that the potential $V(r)$ of a
uniformly charged sphere of radius $R$ and total charge $Ze$ is
\begin{equation*}
V(r)=
\begin{cases}
\dfrac{Ze}{4\pi\epsilon_0 r}, & r\ge R,\\[3mm]
\dfrac{Ze}{4\pi\epsilon_0 R}\left(\dfrac{3}{2}-\dfrac{r^2}{2R^2}\right), & r\le R.
\end{cases}
\end{equation*}

Since the mass of the lead nucleus is
$M_{\mathrm{Pb}}\simeq\dfrac{208\times 1836\,m_e}{207\,m_e}=1.8\times10^3 m_\mu$, the motion of
the lead nucleus may be neglected. When $r>R$, the expression for the Bohr orbit radius of the
$\mu^-$ meson is
\begin{equation}
r_n=\frac{m_e a_0}{m_\mu}\frac{n^2}{Z}=\frac{0.529}{207\times 82}n^2=3.12\,n^2\ (\mathrm{fm}),
\tag{1}
\end{equation}
and the expression for the energy levels is
\begin{equation}
E_n=-\frac{1}{2}m_\mu(\alpha c)^2\frac{Z^2}{n^2}
=-\frac{13.6\times 207\times 82^2}{n^2}=-18.93\,\frac{1}{n^2}\ (\mathrm{MeV}).
\tag{2}
\end{equation}

When $r<R$, the $\mu^-$ meson moves in a circle inside the nucleus; from the centripetal force
formula
\begin{equation}
m_\mu\frac{v^2}{r}=-e\frac{\dd V(r)}{\dd r}=\frac{Ze^2}{4\pi\epsilon_0 R^3}r
\tag{3}
\end{equation}
the kinetic energy of the $\mu^-$ meson is found to be
\begin{equation*}
\frac{1}{2}m_\mu v^2=\frac{Ze^2r^2}{4\pi\epsilon_0\,2R^3}.
\end{equation*}
The potential energy of the $\mu^-$ meson inside the nucleus is
\begin{equation*}
-eV(r)=-\frac{Ze^2}{4\pi\epsilon_0 R}\left(\frac{3}{2}-\frac{r^2}{2R^2}\right).
\end{equation*}
Hence the total energy of the lead $\mu^-$ meson atom is
\begin{align}
E&=\frac{1}{2}m_\mu v^2-eV(r)\notag\\
&=\frac{Ze^2}{4\pi\epsilon_0 R}\left(\frac{r^2}{R^2}-\frac{3}{2}\right). \tag{4}
\end{align}

From the angular momentum quantization condition $m_\mu v r=n\hbar$, $n=1,2,3,\dots$, combined
with (3), the expression for the Bohr orbit radius of the $\mu^-$ meson inside the nucleus is
solved to be
\begin{equation}
r_n=\left(\frac{4\pi\epsilon_0 R^3\hbar^2}{Ze^2m_\mu}n^2\right)^{1/4}
=\left(\frac{137\times 7.7^3\times 197.3}{82\times 207\times 0.511}n^2\right)^{1/4}
=6.14\sqrt{n}\ (\mathrm{fm}). \tag{5}
\end{equation}

Substituting into (4) gives the expression for the energy levels:
\begin{align}
E_n&=\frac{Ze^2}{4\pi\epsilon_0 R}\left(n\hbar\sqrt{\frac{4\pi\epsilon_0}{Ze^2m_\mu R}}-\frac{3}{2}\right)\notag\\[1mm]
&=\frac{82\times 197.3}{137\times 7.7}\left(\frac{6.14^2}{7.7^2}n-\frac{3}{2}\right)\notag\\[1mm]
&=(9.75n-23.00)\ (\mathrm{MeV}). \tag{6}
\end{align}

(1) From (5) and (1) the first two Bohr orbit radii of the $\mu^-$ meson are found to be
\begin{equation*}
r_1=6.14\ \mathrm{fm}\quad(\text{inside the nucleus})
\end{equation*}
and
\begin{equation*}
r_2=12.48\ \mathrm{fm}\quad(\text{outside the nucleus}).
\end{equation*}

(2) From (6) and (2) the energies of the first two levels of the lead $\mu^-$ meson atom are
found to be
\begin{equation*}
E_1=-13.25\ \mathrm{MeV}
\end{equation*}
and
\begin{equation*}
E_2=-\frac{18.93}{4}=-4.73\ (\mathrm{MeV}).
\end{equation*}

(3) The wavelength of the first line of the Lyman series of the lead $\mu^-$ meson atom is
\begin{equation*}
\lambda=\frac{hc}{E_2-E_1}=\frac{1240\times10^{-6}}{13.25-4.73}=1.46\times10^{-4}\ (\mathrm{nm}).
\end{equation*}

(4) If the lead nucleus is regarded as a point charge, then from (2) the energies of the first
two levels of the lead $\mu^-$ meson atom are found to be
\begin{equation*}
E_1'=-18.93\ \mathrm{MeV}
\end{equation*}
and
\begin{equation*}
E_2'=-\frac{18.93}{4}=-4.73\ (\mathrm{MeV}).
\end{equation*}
Hence, in the point-nucleus model, the wavelength of the first line of the Lyman series of the
lead $\mu^-$ meson atom is
\begin{equation*}
\lambda'=\frac{hc}{E_2'-E_1'}=\frac{12400\times10^{-6}}{18.93-4.73}=8.73\times10^{-5}\ (\mathrm{nm}).
\end{equation*}

\begin{commentary}
(1) In using Bohr theory to calculate the energy levels of an atom (or ion), we usually regard
the nucleus as a point charge. In reality every nucleus has a finite size, so only at the
surface of the nucleus or outside it is the potential of the nucleus that of a point charge.
This problem uses Bohr's circular-orbit model to estimate the influence of the finite size of
the nucleus on the energy levels and spectra of exotic atoms. From the results of parts (3)
and (4) it can be seen that, since in the ground state the $\mu^-$ meson has entered the
nucleus so that the effective nuclear charge it experiences is $Z^*e<Ze$, the ground-state
level is raised, and hence $\lambda>\lambda'$. The correction to the ground-state energy of the
atom is
about $30\%$; the correction to the wavelength of the first line of the Lyman series is about
$67\%$. In fact, although $R>r_1$, the $\mu^-$ meson does not stay inside the nucleus all the
time, so the correction obtained by the method of this problem is on the large side.

(2) When the nuclear radius $R<r_1$, according to the method of this problem the above
correction is zero. According to quantum mechanics, even if $R<r_1$, the electron still has a
certain probability of entering the nucleus, and the correction of the finite nuclear size to
the ground-state energy is not zero. Quantum mechanics uses perturbation
theory to obtain the correction of this effect to the ground-state energy of a one-electron
atom (or ion) as
\begin{equation*}
\delta E=\frac{4}{5}\left(\frac{ZR}{a_0}\right)^2\frac{Z^2e^2}{4\pi\epsilon_0\cdot 2a_0}
=\frac{4}{5}\left(\frac{ZR}{a_0}\right)^2\left|E_1'\right|,
\end{equation*}
where $E_1'$ is the ground-state energy in the point-nucleus model, the nuclear radius is
$R=1.3\times A^{1/3}\ \mathrm{fm}$, and the Bohr radius is $a_0=0.0529\ \mathrm{nm}$.

For the neon ion $\mathrm{Ne}^{9+}$ ($Z=10$, $A=20$), $\delta E=3.6\times10^{-7}|E_1'|$; for the
lead ion $\mathrm{Pb}^{82+}$ ($Z=82$, $A=208$), $\delta E=1.1\times10^{-4}|E_1'|$. It can be
seen that for light elements this effect may be neglected altogether, while for heavy elements,
although the effect is very small, it is nevertheless detectable. For the neon $\mu^-$ meson
atom, since the mass of the $\mu^-$ meson is 207 times the mass of the electron,
$\delta E=0.015|E_1'|$, which is $(207)^2$ times larger than that for the neon ion. For the
above three atoms (ions) the nuclear radius $R$ is smaller than $r_1$.

(3) From the data in comment 2 it can be seen that, since the mass of the $\mu^-$ meson is 207
times the mass of the electron, the $\mu^-$ meson will come very close to the nucleus. In this
way we may regard the $\mu^-$ meson as a detector with which to probe the nucleus. The
so-called finite-size effect produced by the size of the nucleus is all the more clearly
displayed in the level structure of such exotic atoms. Hence, if we measure the energy levels
of $\mu^-$ meson atoms, we can infer the charge distribution in the nucleus and the size of the
nucleus. For example, for the $\mu^-$ meson hydrogen atom, provided the level energies are
known to a certain accuracy, the size of the proton can be deduced accurately. In the past the
method of measuring the size of the proton was electron scattering, in which the size was
obtained by analysing the scattering angle.

As an important constituent of matter, the proton's size is of great importance in physics. Only
independent measurements of the proton's size can genuinely make use of the hydrogen spectrum
for a further test of quantum electrodynamics; they can also serve in the future to test
calculations of lattice quantum chromodynamics.
\end{commentary}
\end{solution}

\begin{solution}{5.11}
(1) The first term on the right-hand side of the formula corresponds to the ground state of the
sodium atom, and the second term to an excited state of the sodium atom; the larger $m$ is, the
higher the excited state. The core of the sodium atom may be regarded as an effective nucleus.
When the sodium atom is in a highly excited state $m$, the valence electron is far from the
nucleus and the other 10 inner electrons, and what it experiences is the Coulomb interaction of
a nucleus completely screened by the inner electrons, so the effective nuclear charge may be
taken as 1. Only in this way does the second term on the right-hand side of the formula have
exactly the same form as the hydrogen spectral term; therefore the formula holds only when $m$
is large.

(2) The ionization energy $E_\infty$ of sodium is
\begin{equation*}
E_\infty=0.378Rhc=0.378\times 13.6=5.14\ (\mathrm{eV}).
\end{equation*}

\begin{commentary}
Now let us use the ionization energy $E_\infty$ of sodium to calculate the effective nuclear
charge $Z^*$ felt by the valence electron when the sodium atom is in the ground state. Similarly
to the hydrogen atom, for the sodium atom
\begin{equation*}
E_\infty=\frac{Z^{*2}Rhc}{n^2},
\end{equation*}
where $n=3$ is the principal quantum number of the valence electron when the sodium atom is in
the ground state. From this we obtain
\begin{equation*}
Z^*=\sqrt{\frac{n^2}{Rhc}E_\infty}=\sqrt{\frac{3^2}{13.6}\times 5.14}=1.84.
\end{equation*}

The degree to which the Coulomb potential of the nucleus is screened (that is, the effective
nuclear charge exerted by the core on the valence electron) depends on the state in which the
valence electron finds itself. In different states the probability of the valence electron
appearing at various places in space is different. As $n$ increases, the distance of the valence
electron from the nucleus also increases, while the probability of the valence electron entering
the core becomes smaller, so the degree of screening of the Coulomb potential of the nucleus by
the inner electrons becomes more and more
pronounced, and the effective nuclear charge felt by the valence electron becomes smaller and
smaller. From the above calculation it can be seen that the valence electron of the sodium atom
in the ground state (in the $3s$ state) has a comparatively large probability of appearing near
the nucleus, and therefore feels a comparatively large effective nuclear charge
($Z^*=1.84$).
\end{commentary}
\end{solution}

\begin{solution}{5.12}
(1) The ground-state energy of the helium ion ($\mathrm{He}^+$) and the first Bohr orbit radius
of the electron are
\begin{equation*}
E_1(e)=-\frac{1}{2}m(\alpha c)^2Z^2=-13.6\times 4=-54.4\ (\mathrm{eV}),
\end{equation*}
\begin{equation*}
a_1(e)=\frac{4\pi\epsilon_0\hbar^2}{me^2Z}=\frac{0.0529}{2}=0.026\ (\mathrm{nm}).
\end{equation*}

(2) Similarly to the helium ion, the ground-state energy of the $\mu^-$ meson helium ion and the
first Bohr orbit radius of the $\mu^-$ meson are
\begin{equation*}
E_1(\mu)=-\frac{1}{2}m_\mu(\alpha c)^2Z^2=-13.6\times 4\times 200=-1.1\times10^4\ (\mathrm{eV}),
\end{equation*}
\begin{equation*}
a_1(\mu)=\frac{4\pi\epsilon_0\hbar^2}{m_\mu e^2Z}=\frac{0.0529}{200\times 2}=1.3\times10^{-4}\ (\mathrm{nm}).
\end{equation*}

(3) In the ground state of the $\mu^-$ meson helium atom both the $\mu^-$ meson and the electron
are in the state $n=1$. Since the probability of the $\mu^-$ meson appearing near the helium
nucleus is far greater than that of the electron, it may be said approximately that the electron
is acted on by a nucleus completely screened by the $\mu^-$ meson, i.e. $Z^*\simeq 1$, whereas the
$\mu^-$ meson is acted on by a completely unscreened nucleus, i.e. $Z^*\simeq 2$. Hence

(a) for $Z^*\simeq 1$, the first Bohr orbit radius of the electron is
\begin{equation*}
a_1=\frac{4\pi\epsilon_0\hbar^2}{me^2Z^*}=a_0=0.053\ (\mathrm{nm});
\end{equation*}

(b) the ground-state energy of the $\mu^-$ meson helium atom is
\begin{equation*}
E=E_1(\mu)+E_1(e)=-1.1\times10^4-13.6=-1.1\times10^4\ (\mathrm{eV});
\end{equation*}

(c) the chemical properties of this atom are most similar to those of hydrogen. Since the
$\mu^-$ meson is very close to the nucleus, is tightly bound to the helium nucleus and is not
easily excited, it acts like an effective nucleus of effective nuclear charge $Z^*\simeq 1$.
Outside this effective nucleus there is one electron, and the chemical properties of an atom are
determined by the number of electrons in its outer shell; therefore the $\mu^-$ meson helium
atom resembles hydrogen rather than helium.

\begin{commentary}
The $\mu^-$ meson atom is a kind of exotic atom. Compared with other exotic atoms, the lifetime
of the $\mu^-$ meson atom is comparatively long. It was discovered in the late 1940s by the
famous Chinese high-energy physicist Zhang Wenyu, academician of the Chinese Academy of
Sciences, and has therefore been called the ``Zhang atom'' and ``Zhang radiation''. The
$\mu^-$ meson atom broke through the Rutherford--Bohr atomic model and opened a new field of
research on exotic atoms.
\end{commentary}
\end{solution}

\begin{solution}{5.13}
Let the mass of the particle be $m$ and the mass of the hydrogen atom be $M$, and let the
relative speeds of the particle and the hydrogen atom before and after the collision be
$v_{\text{rel}}$ and $v'_{\text{rel}}$ respectively. Since the momentum of the system is
conserved, the velocity of the centre of mass remains unchanged throughout; then energy
conservation gives
\begin{equation*}
\frac{1}{2}\mu v_{\text{rel}}^2=\frac{1}{2}\mu v_{\text{rel}}'^2+\Delta E,
\end{equation*}
where $\mu$ is the reduced mass, $\mu=\dfrac{mM}{m+M}$, and $\Delta E$ is the excitation energy
acquired by the hydrogen atom. When $v'_{\text{rel}}=0$, the excitation energy acquired by the
hydrogen atom after the collision is maximum, and from the above equation
\begin{equation}
\Delta E=(\Delta E)_{\max}=\frac{1}{2}\mu v_{\text{rel}}^2. \tag{1}
\end{equation}

(1) The first excitation energy of the hydrogen atom is about $10\ \mathrm{eV}$, while at room
temperature the average kinetic energy of the thermal motion of a hydrogen atom is about
$1/25\ \mathrm{eV}$; hence before the collision the hydrogen atom at room temperature is in the
ground state. The H$\alpha$ line is emitted when the hydrogen atom makes a transition from the
$n=3$ state to the $n=2$ state. According to the above analysis, in order that the hydrogen atom
emit the H$\alpha$ line, the hydrogen atom must be excited from the ground state ($n=1$) to the
$n=3$ state. Therefore the excitation energy acquired by the hydrogen atom after the collision
must satisfy $(\Delta E)_{\max}\ge E_3-E_1$. From (1) it is known that this condition is
\begin{equation}
\frac{1}{2}\mu v_{\text{rel}}^2\ge E_3-E_1. \tag{2}
\end{equation}

Let the speed of the electron before the collision be $v$. Compared with $v$, the hydrogen atom
at room temperature before the collision may be regarded as at rest, so
$v_{\text{rel}}=v$; and since $m\ll M$, $\mu\simeq m$. Therefore the kinetic energy of the
electron before the collision must satisfy
\begin{equation*}
\frac{1}{2}mv^2\ge E_3-E_1=-\frac{13.6}{9}+13.6=12.09\ (\mathrm{eV}),
\end{equation*}
that is, the minimum kinetic energy of the electron should be $12.09\ \mathrm{eV}$.

(2) Since the mass $m$ of the proton and the mass $M$ of the hydrogen atom may be regarded as
equal, $\mu=\dfrac{m}{2}$. From (2) we obtain
\begin{equation*}
\frac{1}{2}\mu v_{\text{rel}}^2=\frac{1}{4}mv^2\ge E_3-E_1.
\end{equation*}
Therefore the minimum kinetic energy $E_k$ of the proton should be
\begin{equation*}
E_k=\frac{1}{2}mv^2=2(E_3-E_1)=24.18\ (\mathrm{eV}).
\end{equation*}

(3) As stated in the problem, $\mu=\dfrac{m}{2}$ and $v_{\text{rel}}=2v$. Substituting into (2)
gives
\begin{equation*}
\frac{1}{2}\mu v_{\text{rel}}^2=mv^2\ge E_3-E_1.
\end{equation*}
Therefore the minimum kinetic energy $E_k$ of the proton should be
\begin{equation*}
E_k=\frac{1}{2}mv^2=\frac{1}{2}(E_3-E_1)=6.05\ (\mathrm{eV}).
\end{equation*}

\begin{commentary}
An important physical concept involved in this problem is the following: in the above inelastic
collision process, how much of the kinetic energy of the particle system can be converted
through the collision into the excitation energy of the atom? According to K\"onig's theorem, the
kinetic energy of a system of particles is equal to the sum of the kinetic energy of the centre
of mass and the kinetic energy of the system relative to the centre of mass. Since this system
is not acted on by external forces, the momentum of the system is conserved, and throughout the
process neither the velocity of the centre of mass nor the kinetic energy of the centre of mass
changes. Therefore only the kinetic energy of the system relative to the centre of mass can
possibly be converted through the collision into the excitation energy of the atom.
\end{commentary}
\end{solution}

\begin{solution}{5.14}
(1) If there were no interaction between the electrons, the ground-state energy of the helium
atom would be
\begin{equation*}
E_1=-2R_{\mathrm{He}}hc\frac{Z^2}{n^2}=-2\times 13.6\times\frac{2^2}{1}=-108.8\ (\mathrm{eV}).
\end{equation*}

(2) Neglecting the interaction between the electrons and the motion of the helium nucleus, the
orbit radius of the electron in the ground state is
\begin{equation*}
r_1=a_0\frac{n^2}{Z}=\frac{a_0}{2},
\end{equation*}
where $a_0$ is the Bohr radius. Suppose that the two electrons, owing to their electrostatic
repulsion, tend to move apart, so that their separation is $2r_1$. Then the Coulomb interaction
between the electrons is
\begin{equation*}
E_c=\frac{e^2}{4\pi\epsilon_0\cdot 2r_1}=\frac{14.4}{0.53}=27.2\ (\mathrm{eV}).
\end{equation*}
Taking the Coulomb interaction between the electrons into account, the ground-state energy of
the helium atom is
\begin{equation*}
E_1'=E_1+E_c=-108.8+27.2=-81.6\ (\mathrm{eV}).
\end{equation*}

(3) Since the $4f$ electron is far from the nucleus while the $1s$ electron is close to it, the
$4f$ electron is in the Coulomb field of a completely screened nucleus, with effective nuclear
charge $Z^*\simeq 1$, whereas the $1s$ electron is in the Coulomb field of an unscreened
nucleus, with effective nuclear charge $Z^*\simeq 2$. Using the hydrogen energy-level formula,
the energy of the helium atom is
\begin{equation*}
E_{1s4f}=-13.6\times\frac{2^2}{1^2}-13.6\times\frac{1^2}{4^2}=-55.25\ (\mathrm{eV}).
\end{equation*}

\begin{commentary}
Experimentally the ground-state energy of the helium atom is about $79\ \mathrm{eV}$. Comparing
this with the answers to parts (1) and (2), it is seen that in discussing the energy-level
structure of a multi-electron atom one must take into account, besides the Coulomb interaction
between the electrons and the nucleus, the Coulomb interaction between the electrons themselves.
It is evident that the energy levels of helium are not the same as those of hydrogen, but when
one electron of the helium atom is excited into a state far from the nucleus (a state with large
$n$ and $l$, for example the $4f$ state), its energy level coincides with the corresponding
level of hydrogen (for example the $4f$ level), differing only by a constant ($-54.4\
\mathrm{eV}$).
\end{commentary}
\end{solution}

\begin{solution}{5.15}
Since the half-life $T_{1/2}({}^{238}\mathrm{U})$ of the parent nucleus is much longer than that
of all the daughter nuclei, and among the daughter nuclei the one with the longest half-life is
${}^{234}\mathrm{U}$, with $T_{1/2}({}^{234}\mathrm{U})=2.45\times10^5$ years, after about
$5\times T_{1/2}({}^{234}\mathrm{U})\approx1.2\times10^6$ years the whole decay series reaches
secular equilibrium, that is, the activities of the parent and of all the daughters are equal:
$\lambda_1N_1=\lambda_2N_2=\cdots$. After this, the amount of the final stable nuclide
continuously increases, while the sum of the numbers of daughter nuclei remains essentially
unchanged.

Suppose that the present time is $t_1$ and the initial time is $t_0$. Then at time $t_1$ the sum
of the residual parent nuclei $N_P(t_1)$, the stable nuclei $N_D(t_1)$ and the total number
$\sum_iN_i(t_1)$ of nuclei of all generations of daughters equals the number $N_P(t_0)$ of
parent nuclei at the initial time, i.e.
\begin{equation*}
N_D(t_1)+N_P(t_1)+\sum_iN_i(t_1)=N_P(t_0).
\end{equation*}
Since $N_P(t_1)=N_P(t_0)\ee^{-\lambda(t_1-t_0)}$, that is, in the above equation
$N_P(t_0)=N_P(t_1)\ee^{\lambda(t_1-t_0)}$. In addition, considering that the geological age is
very long and the time required to reach secular equilibrium is much shorter than it, after
secular equilibrium the sum of the daughter nuclei remains unchanged, so approximately
$N_D(t_1)\gg\sum_iN_i(t_1)$. We therefore obtain
\begin{equation*}
N_D(t_1)=N_P(t_1)\left[\ee^{\lambda(t_1-t_0)}-1\right].
\end{equation*}
That is,
\begin{equation*}
\Delta t=t_1-t_0=\frac{1}{\lambda}\ln\left[1+\frac{N_D(t_1)}{N_P(t_1)}\right].
\end{equation*}
Substituting the known quantities, the age of the Earth is estimated to be
\begin{equation*}
\Delta t=\frac{4.468\times10^9\ \text{yr}}{0.693}\ln\left[1+\frac{1}{2.785}\right]
=1.98\times10^9\ (\text{yr}).
\end{equation*}

\begin{commentary}
Through this problem one may gain a good understanding of the laws of radioactive decay, of the
characteristics of the secular equilibrium of radioactive series, and of their practical
applications.
\end{commentary}
\end{solution}

\begin{solution}{5.16}
(1) The disintegrations per minute of 1 gram of carbon come from the ${}^{14}\mathrm{C}$ contained
in it. Let the activity be $A'$; then
\begin{align*}
A'&=\lambda N({}^{14}\mathrm{C})=\lambda\cdot\eta\cdot\frac{1}{12}N_A\\
&=\frac{0.693}{5730\times 365\times 24\times 60}\times 1.2\times10^{-12}\times\frac{1}{12}
\times 6.02\times10^{23}=14\ \text{decays}/(\text{gram}\cdot\text{min}),
\end{align*}
where $N_A$ is Avogadro's constant.

(2) $A(t)=A(0)\ee^{-\lambda t}$. Now $A(0)=A'=14$ decays/(gram$\cdot$min), so the age of the
ancient biological sample is
\begin{equation*}
t=\frac{1}{\lambda}\ln\frac{A(0)}{A(t)}=\frac{5730\ \text{yr}}{0.693}\ln 1.831\approx 5000\ (\text{yr}).
\end{equation*}

(3) $\delta=1\%$ requires a count $N=10000$. Let $m$ be the required number of grams of the
sample; the count of the ancient biological sample in 1 hour is now required to be 10000, hence
\begin{equation*}
14\times 60\times m\cdot\frac{1}{1.831}=10000.
\end{equation*}
From this we obtain
\begin{equation*}
m=\frac{10000}{459}=22\ (\text{grams}).
\end{equation*}

Let the error in the age determined under this error be $\Delta t$; then from the formula
$t=\dfrac{1}{\lambda}\ln\dfrac{A(0)}{A(t)}$ we obtain
\begin{equation*}
\Delta t=\frac{1}{\lambda}\frac{\Delta A(t)}{A(t)}
=\frac{5730\ \text{yr}}{0.693}\times\frac{1}{100}=83\ (\text{yr}).
\end{equation*}

\begin{commentary}
This is a problem on archaeology with ${}^{14}\mathrm{C}$, and it is of practical significance.
\end{commentary}
\end{solution}

\begin{solution}{5.17}
(1) The $\beta^-$ decay formula is
\begin{equation*}
{}_{Z}^{A}\mathrm{X}\to{}_{Z+1}^{A}\mathrm{Y}+\mathrm{e}^-+\bar{\nu}_e.
\end{equation*}
The decay energy is the energy that can be released in the decay process (that is, the maximum
kinetic energy of the emitted electron). Denoting the nuclear mass by $m$, the decay energy is
\begin{equation*}
E_d(\beta^-)=\left[m_X-(m_Y+m_e)\right]c^2.
\end{equation*}
Considering that what is actually known are atomic masses, then, in the case where the
difference between the electron binding energies of the parent (X) and daughter (Y) atoms is
neglected, we have
\begin{equation*}
E_d(\beta^-)=\left[M(\mathrm{X})-M(\mathrm{Y})\right]c^2,
\end{equation*}
where $M$ denotes the atomic mass.

(2) In this problem
\begin{equation*}
E_d(\beta^-)=\left[M({}^{17}\mathrm{N})-M({}^{17}\mathrm{O}^*)\right]c^2=3.72\ (\mathrm{MeV}),
\end{equation*}
that is,
\begin{equation*}
M({}^{17}\mathrm{O}^*)c^2=M({}^{17}\mathrm{N})c^2-3.72\ (\mathrm{MeV}).
\end{equation*}

If ${}^{17}\mathrm{O}^*$ can emit a neutron and reach the excited state ${}^{16}\mathrm{O}^*$,
then energetically the following inequality must be satisfied:
\begin{equation*}
M({}^{17}\mathrm{O}^*)c^2-\left[M({}^{16}\mathrm{O})c^2+E^*({}^{16}\mathrm{O})\right]-m_nc^2>0,
\end{equation*}
where $E^*({}^{16}\mathrm{O})$ is the excitation energy of ${}^{16}\mathrm{O}$. Substituting the
expression for $M({}^{17}\mathrm{O}^*)c^2$ into the above gives
\begin{align*}
&\left[M({}^{17}\mathrm{N})-M({}^{16}\mathrm{O})-m_n\right]c^2-3.72\ \mathrm{MeV}-E^*({}^{16}\mathrm{O})\\
&\quad=(17.008449-15.994915-1.008665)uc^2-3.72\ \mathrm{MeV}-E^*({}^{16}\mathrm{O})\\
&\quad=0.82\ \mathrm{MeV}-E^*({}^{16}\mathrm{O})>0.
\end{align*}
Since all the excitation energies $E^*({}^{16}\mathrm{O})$ are greater than $0.82\ \mathrm{MeV}$,
in this problem ${}^{17}\mathrm{O}^*$ can emit a neutron in only one way, namely by decaying to
the ground state of ${}^{16}\mathrm{O}$.

Let the energy of the emitted neutron, with the recoil of ${}^{16}\mathrm{O}$ taken into account,
be $E_n$; then
\begin{equation*}
E_n+E_n\cdot\frac{m_n}{M({}^{16}\mathrm{O})}=0.82\ (\mathrm{MeV}).
\end{equation*}
Hence
\begin{equation*}
E_n\approx 0.82\times\frac{16}{17}\ \mathrm{MeV}=0.77\ (\mathrm{MeV}).
\end{equation*}

\begin{commentary}
The decay energy is an important physical quantity in nuclear decay. This problem is concerned
with how to use energy relations to decide whether a particular nuclear reaction can occur.
\end{commentary}
\end{solution}

\begin{solution}{5.18}
(1) The semi-empirical formula for the binding energy of a nucleus with $Z$ protons and $N$
neutrons ($A=Z+N$) is (the famous Weizs\"acker formula)
\begin{equation*}
B(Z,A)=a_V A-a_sA^{2/3}-a_c\frac{Z^2}{A^{1/3}}-a_a\frac{(N-Z)^2}{A}+B_P,
\end{equation*}
where $a_V$, $a_s$, $a_c$, $a_a$ are adjustable parameters (constants), $a_V A$ is the volume
energy, $a_sA^{2/3}$ the surface energy, $a_c\dfrac{Z^2}{A^{1/3}}$ the Coulomb energy,
$a_a\dfrac{(N-Z)^2}{A}$ the symmetry energy, and $B_P$ the pairing energy. The empirical formula
for $B_P$ is $B_P=\delta a_PA^{-1/2}$, where $a_P$ is a constant,
\begin{equation*}
\delta=
\begin{cases}
1 & \text{even--even nucleus},\\
0 & \text{odd--even nucleus},\\
-1 & \text{odd--odd nucleus}.
\end{cases}
\end{equation*}
The last two terms in the formula arise from quantum effects and have no classical counterpart.

(2) In order to use the semi-empirical formula to calculate the nuclear radius constant, the
Coulomb energy term in it should be replaced by the more accurate formula for the Coulomb energy
that contains the nuclear radius $R$:
\begin{equation*}
B_c=\frac{3}{5}\frac{Z(Z-1)e^2}{4\pi\epsilon_0R},
\end{equation*}
in which $R=\gamma_0A^{1/3}$, where $\gamma_0$ is the nuclear radius constant; the corresponding
value of $\gamma_0$ can be determined from the mass difference of mirror nuclei.

Knowing the binding energy of a nucleus, one obtains the corresponding semi-empirical formula for
the atomic mass:
\begin{equation*}
M(Z,A)=ZM({}^{1}\mathrm{H})+Nm_n-\frac{B(Z,A)}{c^2}.
\end{equation*}
Then for the mirror nuclei ${}^{11}_{5}\mathrm{B}$ and ${}^{11}_{6}\mathrm{C}$ we have
\begin{equation*}
M({}^{11}_{5}\mathrm{B})=5M({}^{1}\mathrm{H})+6m_n-B({}^{11}_{5}\mathrm{B})/c^2,
\end{equation*}
\begin{equation*}
M({}^{11}_{6}\mathrm{C})=6M({}^{1}\mathrm{H})+5m_n-B({}^{11}_{6}\mathrm{C})/c^2.
\end{equation*}
That is,
\begin{equation*}
M({}^{11}_{6}\mathrm{C})-M({}^{11}_{5}\mathrm{B})=M({}^{1}\mathrm{H})-m_n
+B({}^{11}_{5}\mathrm{B})/c^2-B({}^{11}_{6}\mathrm{C})/c^2.
\end{equation*}
Since the two mirror nuclei have the same $A$ and the same $(N-Z)^2$, we have
\begin{align*}
\left[M({}^{11}_{6}\mathrm{C})-M({}^{11}_{5}\mathrm{B})\right]c^2
&=\left[M({}^{1}\mathrm{H})-m_n\right]c^2
-\frac{3}{5}(5\times 4)\frac{e^2}{4\pi\epsilon_0\gamma_0A^{1/3}}\\
&\quad+\frac{3}{5}(6\times 5)\frac{e^2}{4\pi\epsilon_0\gamma_0A^{1/3}}.
\end{align*}
That is,
\begin{equation*}
\frac{6e^2}{4\pi\epsilon_0\gamma_0A^{1/3}}
=\left[M({}^{11}_{6}\mathrm{C})-M({}^{11}_{5}\mathrm{B})-M({}^{1}\mathrm{H})+m_n\right]c^2
=2.76\ (\mathrm{MeV}),
\end{equation*}
\begin{equation*}
\gamma_0=\frac{6\times 1.44\ \mathrm{MeV}\cdot\mathrm{fm}}{2.76\ \mathrm{MeV}\times 11^{1/3}}
=1.41\ (\mathrm{fm}).
\end{equation*}

\begin{commentary}
The liquid-drop model description of the nucleus is very important in nuclear physics. Through
this problem the reader is required to become familiar with the expression of the semi-empirical
binding-energy formula and with its simple applications.
\end{commentary}
\end{solution}

\begin{solution}{5.19}
In this problem one may introduce the representation of the physical quantities in the
centre-of-mass frame; combining the centre-of-mass frame with the laboratory frame makes the
discussion convenient and clear. As shown in the figure, the incident particle $M_1$ is
$M({}^{4}_{2}\mathrm{He})$ and the target nucleus $M_2$ is $M({}^{1}_{1}\mathrm{H})$. The speed of
the incident particle is $v_1$, and the target nucleus is assumed to be at rest. The speed of the
outgoing particle in the $L$ frame is $v_1'$, and in the centre-of-mass frame it is $V_1$; the
recoil speed of the target nucleus in the $L$ frame is $v_2'$, and in the centre-of-mass frame
it is $V_2$.

\begin{figbox}{5.19}
\begin{tikzpicture}[scale=1.1]
  % collision point and CM construction
  \coordinate (A) at (0,0);
  \coordinate (C) at (2.55,0);
  \coordinate (B) at (4.5,-1.62);
  \coordinate (O) at (2.05,2.65);
  \draw[dashedline] (-1.6,0) -- (4.9,0);
  % incident particle
  \draw[force] (-1.55,0) -- (-0.35,0) node[midway,above] {$v_1$};
  \fill (A) circle (1.5pt);
  \node[below left,inner sep=2pt] at (A) {$A$};
  \node[above=13pt,inner sep=1pt] at (A) {$M_2$};
  \node[above] at (-1.75,0) {$M_1$};
  % outgoing particle
  \draw[force] (A) -- (O) node[midway,above left,inner sep=1pt] {$v_1'$};
  \fill (O) circle (1.5pt);
  \node[above right,inner sep=1pt] at (O) {$O$};
  % CM velocity and V1
  \draw[axis] (A) -- (C) node[midway,above=1pt] {$v_C$};
  \draw[force] (C) -- (O) node[midway,right,inner sep=1pt] {$V_1$};
  \fill (C) circle (1.5pt);
  \node[above right,inner sep=1pt] at (C) {$C$};
  % recoil
  \draw[force] (A) -- (B) node[midway,below right,inner sep=1pt] {$v_2'$};
  \draw[force] (C) -- (B) node[midway,below,inner sep=1pt] {$V_2$};
  \fill (B) circle (1.5pt);
  \node[right] at (B) {$B$};
  % angle at A
  \draw (1.0,0) arc (0:-19.8:1.0);
  \node at (1.18,-0.40) {$\varphi_L$};
\end{tikzpicture}
\end{figbox}

(1) Before the collision the speed of the centre of mass is
$v_C=v_1\dfrac{M_1}{M_1+M_2}$, in the same direction as $v_1$. The speed of the target nucleus is
also $v_C$, but its direction is opposite to that of $v_1$.

Since in the elastic scattering process the magnitudes of the centre-of-mass velocities of the two
colliding particles remain unchanged, the speed of the recoiling target nucleus is $V_2=v_C$;
hence from the figure it is seen that triangle $ABC$ is isosceles. Since the velocity of the
recoiling target nucleus in the laboratory frame is $v_2'=v_C+V_2$, the magnitude of $v_2'$ is
\begin{equation*}
v_2'=2v_C\cos\varphi_L.
\end{equation*}
Accordingly the kinetic energy of the recoiling hydrogen nucleus is
\begin{align*}
E_{\text{recoil}}&=\frac{1}{2}M_2v_2'^2=\frac{1}{2}M_2(2v_C\cos\varphi_L)^2\\
&=\frac{1}{2}M({}^{1}\mathrm{H})\left[\frac{2M({}^{4}_{2}\mathrm{He})v_1}
{M({}^{1}\mathrm{H})+M({}^{4}_{2}\mathrm{He})}\right]^2\cos^2 30\degree\\
&=\frac{4M({}^{1}\mathrm{H})M({}^{4}_{2}\mathrm{He})}
{\left[M({}^{1}\mathrm{H})+M({}^{4}_{2}\mathrm{He})\right]^2}E({}^{4}_{2}\mathrm{He})\cos^2 30\degree\\
&\approx\frac{4\times 1\times 4}{(1+4)^2}\cdot 2\ \mathrm{MeV}\cdot\left(\frac{\sqrt{3}}{2}\right)^2
=0.96\ (\mathrm{MeV}).
\end{align*}

(2) The first step is to find the expression for the differential cross section of the recoil
particle in the centre-of-mass frame, $\sigma_C(\varphi_C)$. Using
$\sigma_C(\varphi_C)\dd\Omega(\varphi_C)=-\sigma_C(\theta_C)\dd\Omega(\theta_C)$ (the minus sign
arises because when $\dd\Omega(\varphi_C)$ increases, the corresponding $\dd\Omega(\theta_C)$
decreases), we obtain
\begin{equation*}
\sigma_C(\varphi_C)=-\frac{\sin\theta_C\dd\theta_C}{\sin\varphi_C\dd\varphi_C}\sigma_C(\theta_C).
\end{equation*}
Since $\theta_C=\pi-\varphi_C$, we have $\dd\theta_C=-\dd\varphi_C$ and $\sin\theta_C=\sin\varphi_C$,
and hence
\begin{equation*}
\sigma_C(\varphi_C)=\left(\frac{a}{4}\right)^2\frac{1}{\sin^4\left(\dfrac{\pi-\varphi_C}{2}\right)}
=\left(\frac{a}{4}\right)^2\frac{1}{\cos^4\left(\dfrac{\varphi_C}{2}\right)}.
\end{equation*}

The second step is to find the expression for the differential cross section of the recoil
particle in the laboratory frame, $\sigma_L(\varphi_L)$. Using
$\sigma_L(\varphi_L)\dd\Omega(\varphi_L)=\sigma_C(\varphi_C)\dd\Omega(\varphi_C)$, we obtain
\begin{equation*}
\sigma_L(\varphi_L)=\frac{\sin\varphi_C\dd\varphi_C}{\sin\varphi_L\dd\varphi_L}\sigma_C(\varphi_C).
\end{equation*}
Since $\varphi_C=2\varphi_L$, we have
$\sin\varphi_C\dd\varphi_C=4\sin\varphi_L\cos\varphi_L\dd\varphi_L$, and hence
\begin{align*}
\sigma_L(\varphi_L)&=4\cos\varphi_L\,\sigma_C(\varphi_C)
=4\cos\varphi_L\left(\frac{a}{4}\right)^2\frac{1}{\cos^4\varphi_L}
=\left(\frac{a}{4}\right)^2\frac{4}{\cos^3\varphi_L}.
\end{align*}

Now
\begin{equation*}
a=\frac{Z_1Z_2e^2}{4\pi\epsilon_0E_C}
=\frac{2\times 1\times 1.44\ \mathrm{eV}\cdot\mathrm{nm}}
{2\times10^6\ \mathrm{eV}\cdot\dfrac{1}{1+4}}=7.2\times10^{-15}\ (\mathrm{m})
\end{equation*}
(the above has made use of the centre-of-mass energy of the incident particle
$E_C=\dfrac{M_2}{M_1+M_2}E_L$).

Therefore
\begin{align*}
\sigma_L(\varphi_L)&=\left(\frac{7.2\times10^{-15}}{4}\right)^2\cdot
\frac{4}{\left(\dfrac{\sqrt{3}}{2}\right)^3}\ (\mathrm{m}^2/\mathrm{sr})\\
&=0.20\ \mathrm{b/sr}\qquad (1\ \mathrm{b}\ (\text{barn})=10^{-28}\ \mathrm{m}^2).
\end{align*}

\begin{commentary}
This problem involves a comprehensive calculation of the kinematics and dynamics of two-particle
elastic scattering, and also involves the transformation of the differential cross section
between the centre-of-mass frame and the laboratory frame.
\end{commentary}
\end{solution}

